\documentclass[journal]{IEEEtran}

\usepackage{amsmath,amssymb}
\usepackage{booktabs}
\usepackage{graphicx}
\usepackage{hyperref}
\usepackage{xcolor}
\usepackage{placeins}
\usepackage{balance}
\hypersetup{hidelinks}

\definecolor{draftred}{RGB}{150,20,20}
\newcommand{\draftwarning}{\textcolor{draftred}{DRAFT---NOT FOR SUBMISSION}}
\newcommand{\TopoCount}{TBD}
\newcommand{\TopoRate}{TBD}
\newcommand{\GeomCount}{TBD}
\newcommand{\GeomRate}{TBD}
\newcommand{\ACTPointerCount}{TBD}
\newcommand{\ACTRelCount}{TBD}
\newcommand{\TeacherCount}{TBD}
\newcommand{\MaxHolmP}{TBD}
\newcommand{\MuJoCoTopoCount}{TBD}
\newcommand{\MuJoCoGeomCount}{TBD}
\IfFileExists{../../artifacts/papers/branchbot/final/counterfactual/branch_numbers.tex}
  {\input{../../artifacts/papers/branchbot/final/counterfactual/branch_numbers.tex}}{}

\title{Semantic Counterfactuals Expose Branch-Identity Shortcuts in
Bimanual Wire-Harness Manipulation}
\author{Anonymous Authors}
\date{}

\begin{document}
\maketitle

\begin{center}
\draftwarning
\end{center}

\begin{abstract}
Bimanual manipulation of a branched wire harness must preserve endpoint
identity as well as geometry. We introduce an exact semantic counterfactual in
which cable state, dynamics, fixtures, and random seed are fixed while the
observable endpoint--arm--target binding is renamed. Across 600
physical-seed--assignment episodes per policy, the full relational model
succeeds in \TopoCount{} (\TopoRate\%), typed ACT variants in
\ACTPointerCount{} and \ACTRelCount, and the teacher in \TeacherCount. Geometry
succeeds in \GeomCount{} (\GeomRate\%); every success occurs under the canonical
assignment and none after the A/B swap. Geometry is lower in all six paired
cells after Holm correction. Full and ablated semantic models remain tied near
the teacher ceiling, showing that the benefit is not unique to message passing:
it follows from preserving observable identity through representation and
typed action grounding. A direct current-pose MuJoCo backend yields
\MuJoCoTopoCount/40 successes for the relational model and
\MuJoCoGeomCount/40 for Geometry; all Geometry successes again occur only under
the canonical assignment. The evidence establishes a semantic-shortcut failure
under a fixed two-branch abstraction. It does not establish 3--5 branch
generalization, image-based identity perception, force control, or production
performance.
\end{abstract}

\begin{IEEEkeywords}
Bimanual manipulation, deformable object manipulation, graph representations,
semantic grounding, wire-harness assembly.
\end{IEEEkeywords}

\section{Introduction}

Wire-harness automation must preserve identity while controlling shape.  Two
branches can be nearly identical in length and geometry yet terminate in
different connectors or vehicle subsystems.  Swapping them may leave a tidy,
crossing-free harness that is electrically and procedurally wrong.  A useful
controller must therefore answer a relational question---which physical
endpoint belongs to which arm and destination---before executing geometric
motion.

Prior systems provide strong real-robot evidence for dense-knot recovery
\cite{sundaresan2021untangling}, bilateral long-cable untangling
\cite{viswanath2022long}, learned tracing and manipulation
\cite{viswanath2023handloom}, and dual-arm DLO planning in constrained 3-D
environments \cite{yu2022coarse,yu2023wholebody}.  BranchBot does not claim to
be the first bimanual cable system.  We instead isolate a narrower failure
mode that these systems do not directly measure: sensitivity to a pure
renaming of branch identity when all physical evidence is held fixed.

Our contributions are:

\begin{enumerate}
  \item a paired A/B intervention that exactly reuses each physical seed while
  changing only observable endpoint--arm--target relations;
  \item a matched family of graph, Transformer, and typed-pointer controllers
  trained on balanced, episode-disjoint demonstrations, revealing which
  semantic channels are sufficient rather than crediting one architecture;
  \item exact paired statistics, semantic-invariance diagnostics, a mutually
  exclusive failure taxonomy, and direct closed-loop MuJoCo replication;
  \item an explicit evidence boundary separating two-branch semantic
  assignment from multi-branch, vision, force, and hardware claims.
\end{enumerate}

\section{Related Work}

\paragraph{Cable untangling and tracing.}
Learned grasp features and recovery policies untangle dense non-planar knots
on real cables \cite{sundaresan2021untangling}; SGTM combines RGB-D perception
with bilateral motion primitives for long cables \cite{viswanath2022long}.
HANDLOOM learns traces and crossing classes from simulated and real RGB-D data
and supports bimanual manipulation \cite{viswanath2023handloom}.  Those works
provide stronger perception and hardware evidence than ours.  Their public
data lack matched branched endpoint--target renamings, so we do not report them
as real-world semantic completion trials.

\paragraph{Dual-arm DLO control.}
Coarse-to-fine and whole-body frameworks combine global planning with
closed-loop control under obstacles and model uncertainty
\cite{yu2022coarse,yu2023wholebody}.  Tactile or force-aware industrial
harnessing further demonstrates why task-space mocap cannot stand in for
contact regulation \cite{zhang2024twisting}.  Our fixed two-branch simulator
is designed for many exact counterfactual pairs rather than physical fidelity.

\paragraph{Learning representations and temporal actions.}
Graph networks have modeled cable deformation \cite{wang2022offlineonline};
graph structure alone is not novel.  ACT predicts temporally chunked bimanual
actions \cite{zhao2023act}, while role assignment can simplify bimanual
coordination \cite{grannen2023stabilize}.  We adapt ACT to the same typed state
and two-pointer action interface so a continuous-output mismatch cannot create
an artificial advantage for the graph policy.

\section{Problem Formulation}

Let $G_t=(V_t,E_t^p,E_t^s)$ contain observable cable particles, target and arm
nodes.  $E_t^p$ encodes cable connectivity; $E_t^s$ connects each semantic
endpoint to its arm, two clearance waypoints, and destination.  A rollout is
successful only if both semantic endpoint errors fall within the difficulty-
dependent tolerance, projected cable crossings are zero, and the condition is
stable for four steps.

For a frozen physical seed $s$, intervention $z\in\{0,1\}$ either retains or
swaps the A/B names.  Cable coordinates, target coordinates, initial arms,
noise, material parameters, and integration are identical.  When $z$ changes,
the endpoint--arm--target relation and corresponding action slots change
together.  A correct equivariant policy should preserve success, while a
fixed-name or position shortcut can silently manipulate the wrong endpoint.

\section{Method}

The PBD harness \cite{muller2007pbd} contains a ten-particle trunk and two
nine-particle branches joined at one junction.  Each node stores normalized
position, generic type, endpoint/junction role, A/B identity, process order,
state, progress, and difficulty.  TopoHarness performs three message-passing
layers over $\mathrm{clip}(E^p+E^s,0,1)$ and produces two node-pointer heads in
addition to six normalized action logits.  Pointer slot A is masked to
observable A targets and slot B to B targets.  Selected coordinates replace
the continuous position outputs; both arms retain their assigned endpoints.

The Geometry model has the same hidden width and training schedule but removes
adjacency, endpoint/junction identity, process order, and A/B features.  Three
graph ablations retain only physical edges (leaving unary identity), retain
only semantic edges, or retain both relations while removing unary ID/order
features.  The ACT-style baseline is an independent compact implementation
with a CVAE Transformer, eight-action chunks, temporal ensembling, and the same
data.  ACT-Pointer adds legal A/B candidate sets at action grounding.
ACT-RelPool additionally averages fixed one-hop semantic-neighbor features
before its Transformer; it does not use learned graph propagation.  Thus the
suite tests several distinct ways of carrying semantic identity to the action,
not merely graph versus non-graph architecture.

\section{Experimental Protocol}

The raw corpus contains 420 complete teacher rollouts and 43,258 sampled
states.  Because the teacher can fail on difficult cable states, model fitting
uses a frozen successful-only subset of 160/40/80 training/validation/test
episodes (25,740 states; SHA-256 prefix \texttt{2e52208d2f41}).  Every split is
balanced across difficulty and A/B assignment; whole episodes and physical
seeds are disjoint.  The final evaluation starts at seed 66,000,000, outside
all corpus seed ranges, and contains 100 new physical seeds at each difficulty
$0.2,0.5,0.8$.  Both semantic assignments are executed for every seed.

All learned graph policies use 64 hidden units and three layers.  ACT variants
use 48 hidden units, one encoder/decoder layer, and eight-step chunks.  Models
see identical demonstrations within their family; model selection uses only
validation loss.  We report Wilson intervals, per-difficulty and per-assignment
success, endpoint errors, crossings, motion length, failures, and paired exact
McNemar tests with Holm family-wise correction.  A single training seed per
model remains a limitation; paired evaluation seeds quantify rollout, not
training, uncertainty.

\begin{figure}[t]
  \centering
  \includegraphics[width=\columnwidth]{%
    ../../artifacts/mujoco_validation/branch_mujoco_final.png}
  \caption{BranchBot visual digital twin: two UR10e arms and Robotiq 2F-85
  grippers around a native MuJoCo trunk--junction--branch cable.  Robot meshes
  are collision-disabled visual IK followers and do not alter cable physics.}
  \label{fig:visual-twin}
\end{figure}

\section{Results}

\begin{table}[t]
  \centering
  \caption{Frozen PBD counterfactual.  Each policy receives 300 physical seeds
  under both semantic assignments (600 episodes).}
  \label{tab:main}
  \IfFileExists{../../artifacts/papers/branchbot/final/counterfactual/branch_paper_table.tex}
    {{\small\setlength{\tabcolsep}{2.5pt}%
      \input{../../artifacts/papers/branchbot/final/counterfactual/branch_paper_table.tex}}}
    {\fbox{\parbox{\dimexpr\columnwidth-2\fboxsep-2\fboxrule\relax}{%
      \centering Frozen table generated by \texttt{sim-eval-branch-paper}}}}
\end{table}

\begin{figure}[t]
  \centering
  \IfFileExists{../../artifacts/papers/branchbot/final/counterfactual/branch_counterfactual_heatmap.pdf}
    {\includegraphics[width=\columnwidth]{%
      ../../artifacts/papers/branchbot/final/counterfactual/branch_counterfactual_heatmap.pdf}}
    {\fbox{\parbox[c][1.25in][c]{0.94\columnwidth}{\centering
      Frozen difficulty--assignment heatmap}}}
  \caption{Success under the exact A/B intervention (100 paired physical seeds
  per cell).  Both columns share identical physics; only semantic binding is
  renamed.}
  \label{fig:heatmap}
\end{figure}

\paragraph{Paired semantic intervention.}
The frozen main table contains \TopoCount{} relational successes and
\GeomCount{} Geometry successes.  TopoHarness differs on only 2 of 300
canonical/swapped physical pairs (99.3\% outcome agreement).  Geometry obtains
all \GeomCount{} successes in the canonical condition and 0/300 after the A/B
swap; its paired outcome agreement is 59.0\%.  Geometry is lower than
TopoHarness in every difficulty--assignment cell after Holm correction, with
the largest adjusted value \MaxHolmP.  Untyped ACT reaches 432/600 and exhibits
15 canonical-success/swapped-failure pairs versus one in the reverse
direction.  At difficulty 0.8 it is lower than TopoHarness for both assignments
(adjusted $p=5.94\!\times\!10^{-3}$ and $4.81\!\times\!10^{-6}$).  ACT-Pointer
reaches \ACTPointerCount{} and ACT-RelPool \ACTRelCount, both near the scripted
teacher's \TeacherCount.  The exact paired tests avoid treating the two
semantic assignments of one physical seed as independent.

\paragraph{Ablations and interpretation.}
Removing semantic edges (while retaining unary identity), removing physical
edges (while retaining semantic identity), or removing unary ID/order (while
retaining semantic relations) yields 477, 478, and 481 successes, respectively,
versus 480 for TopoHarness.  Typed ACT also reaches the same band.  The data
therefore reject a unique message-passing explanation: multiple semantic
channels are sufficient when the two simultaneous commands remain type-bound.
The unsupported alternative is geometry alone, which learns a canonical
left/right shortcut and collapses under renaming.

\begin{table}[t]
  \centering
  \caption{Direct current-pose MuJoCo observation/task-space-control audit.
  The detailed arms are visual IK twins, not force controllers.}
  \label{tab:mujoco}
  \IfFileExists{../../artifacts/papers/branchbot/final/mujoco_direct/branch_mujoco_direct_table.tex}
    {{\small\setlength{\tabcolsep}{2.5pt}%
      \input{../../artifacts/papers/branchbot/final/mujoco_direct/branch_mujoco_direct_table.tex}}}
    {\fbox{\parbox{\dimexpr\columnwidth-2\fboxsep-2\fboxrule\relax}{%
      \centering Frozen direct-MuJoCo table}}}
\end{table}

\paragraph{Cross-physics closed loop.}
For each of 20 paired fixture and initial-gripper perturbations, the graph is
rebuilt from current MuJoCo cable bodies, sites, and mocap poses at every
control step.  Both policy outputs directly move the two task-space grippers;
no PBD trajectory, stage event, or success label is replayed.  TopoHarness
completes \MuJoCoTopoCount/40 semantic trials and Geometry completes
\MuJoCoGeomCount/40.  TopoHarness is 20/20 under each assignment; all 13
Geometry successes are canonical and its swapped score is 0/20.  This
replicates the semantic-shortcut diagnosis under a second deformable backend
\cite{todorov2012mujoco}, but does not validate joint torque or contact-force
control.

\begin{figure}[t]
  \centering
  \IfFileExists{../../artifacts/papers/branchbot/final/mujoco_direct/branch_mujoco_direct_topology_final.png}
    {\includegraphics[width=\columnwidth]{%
      ../../artifacts/papers/branchbot/final/mujoco_direct/branch_mujoco_direct_topology_final.png}}
    {\includegraphics[width=\columnwidth]{%
      ../../artifacts/mujoco_validation/branch_mujoco_final.png}}
  \caption{Representative direct-observation MuJoCo final state.  Native cable
  endpoint--fixture distances define success; the robot mesh only visualizes
  task-space commands.}
  \label{fig:mujoco-final}
\end{figure}
\FloatBarrier

\section{Limitations and Claim Boundary}

The central intervention contains exactly two branches and two labels.  It
does not establish permutation-equivariant control for three to five branches,
sequential regrasp planning, or non-symmetric branch materials.  PBD is planar;
its crossing count is a projection metric.  MuJoCo uses 3-D elastic cables but
task-space mocap grippers and collision-disabled visual robots.  No experiment
uses raw camera images, detects branch identity under occlusion, measures
contact force, or runs on hardware.  HANDLOOM is relevant public perception
data but lacks the matched branched semantic labels required by our causal
test, so it is not counted as real-world completion evidence.  These gaps keep
the manuscript a research draft.

\section{Ethical and Societal Considerations}

The study uses synthetic robot trajectories and no personal data.  Harness
automation could reduce repetitive ergonomic strain but may also displace
manual assembly work.  The simulated controller is not safety certified and
must not command a production cell without guarding, risk assessment, force
limits, and qualified human supervision.  We therefore expose unsuccessful
episodes and avoid translating simulator rates into labor or financial claims.

\section{Reproducibility}

The package includes both source and successful-only datasets, whole-episode
split manifests, hashes, frozen checkpoints, exact seeds, per-episode CSV/JSON,
statistical audit scripts, plotting code, MJCF, visual-asset licenses, and a
one-command build.  The ACT baseline is an independent adaptation and is not
presented as upstream ACT code.  An audit recomputes counts, Wilson inputs,
semantic pairs, exact tests, checkpoint hashes, and PBD/MuJoCo seed separation
before the PDF is copied to the release directory.

\balance
\bibliographystyle{plain}
\bibliography{references}

\end{document}
